\documentclass{article}
\pagenumbering{gobble}
\usepackage{enumerate}
\usepackage[margin=1.4in]{geometry}
\title{Statement on Diversity, Equity and Inclusion}
\author{Joel Villatoro}
\date{}

\usepackage{blindtext}
\usepackage{hyperref}
\begin{document}
\vspace{-0.5in}
\maketitle

Due to my mixed Mexican-American heritage, I have always seen promoting equity and diversity as a vital component of my work. I am deeply committed to promoting an environment that is welcoming to people of all backgrounds. At your institution, I will commit to furthering its mission to provide an equitable and inclusive environment for its faculty and its students.

The lack of diversity in mathematics is just one symptom of many systemic issues that make it more difficult for people from underrepresented groups to succeed. Therefore, in order to promote equity, I believe that it is important to take proactive measures to promote diversity as well. Let me describe some of the activities I have participated in along these lines.

While I was an undergraduate student, I worked as a camp counselor for a day camp that provided critical childcare and education services for poor, largely latino, neighborhoods in Oklahoma City. While a graduate student at the University of Illinois, I participated in a mentorship program for first year graduate students from underrepresented backgrounds. As an MPS-Ascend postdoc, promoting diversity was an important component of my grant application. While at Washington University I have done outreach for high schoolers from predominantly latino communities and served on the mathematics department's DEI committee. I have also been regularly attending the Field of Dreams Conference to support the university's efforts to recruit URM undergraduates. Finally, I helped to restart the undergraduate math club which had become inactive since COVID. One of my primary motivations for this was to place myself in a role where I could steer the undergraduate mathematics community in a direction that was inclusive to people of all backgrounds.

Most of these activities I just described are focused on one-on-one mentorship and contact with people from underrepresented groups. This is due to the fact that I believe that these kinds of mentorship activities are the best way to promote educational success for people from underrepresented groups.

My personal story is one that illustrates the critical importance of having role models. My mother moved to the United States from Mexico when I was five years old. She worked incredibly hard to provide for her four children as a single mother and we frequently had trouble making ends meet. However, while I was growing up, my mother was attending graduate school to get her PhD. Her pay as a graduate student was very low and it was not easy for her to support four children while also pursuing her education. 

Although our financial circumstances were poor, I was able to go to college and eventually get a PhD myself. I attribute this to the fact that I had a readily accessible role model that instilled in me the idea that reaching a high level of education was something that I could achieve. In my work as an academic, I consider it an important responsibility of mine to provide mentorship and guidance to people from diverse backgrounds. My intention is to instill in them the belief that they are capable of joining our academic community and that we welcome them.

\end{document}
